% P50 | 5 Part C summary | What Part C shows about the American put
% guide: dev_notes/paragraph_guides/P50_what_part_c_shows_about_the_american_put.txt
The CRR tree, a finite-difference solver for the variational inequality and least-squares Monte Carlo agree with each other within their numerical errors, and the tree approaches the closed-form perpetual put as the maturity and the number of steps grow.
The one row that disagreed, $T=1000$ with $800$ steps, was a lattice artefact, and the refined computation says so.
Early exercise adds a few percent to an at-the-money put and much more deep in the money, and it is zero for $r=0$ and for calls without dividends.
For the NVIDIA put it adds about \NvdaPutEEPpct\% to the European price.
The buyer's expected profit under the real-world measure depends on the drift, which the price does not.
A wrong volatility corrupts the exercise rule asymmetrically, because exercising too early costs much more than exercising too late.
In practice the tree is the natural tool for prices, boundaries and hedges, a PDE solver is better when accuracy is needed in one or two dimensions, regression Monte Carlo takes over when the dimension is high, and Barone--Adesi--Whaley serves for a quick estimate.
